%%
%% OBJECTS/TIME
%%
\section{Le temps:~les cycles}

Le temps~$t$ est un ensemble de~$N$ espaces ordonn\'es. %
Il supporte toutes les op\'erations sur les s\'equences d'espaces. %
Ces op\'erations sont plus difficiles \`a d\'efinir que celles portant sur
les espaces. %
Leur mise en oeuvre n\'ecessite l'ench\^assement du temps dans un
s\'equenceur. %
C'est pr\'ecis\'ement le choix des primitives appartenant en propre au
temps, par opposition \`a celles qui doivent \^etre associ\'ees \`a un
s\'equenceur qui pose probl\`eme. %
Nous esp\'erons que ces probl\`emes seront peu \`a peu lev\'es \`a l'usage. %
Nous tenterons n\'eanmoins de d\'efinir, en faisant momentan\'ement
abstraction de la structure particuli\`ere d'un s\'equenceur, une liste
d'op\'erations possibles. %

Le temps est a priori une s\'equence d'espaces g\'en\'er\'ee par l'application
it\'erative d'une fonction de transition sur un espace initial. %
Si l'on utilise une m\'etaphore cin\'ematographique, il permet
l'enregistrement, le tirage et le montage de films. %
Il permet \'egalement des r\'eductions d'informations diff\'erentielles sur
les espaces, des r\'eductions de s\'equences de r\'eductions d'espaces, et
des transformations arbitraires. %

\myplotfullR[t]{.75}{anneal-sum}%
{Visualisation par surface d'\'el\'evation d'une r\'eduction int\'egrale du temps}%
{Somme de 400~g\'en\'erations de la r\`egle~\arule{Anneal}}%
{Cette visualisation par surface d'\'el\'evation d'une r\'eduction int\'egrale
  du temps r\'esultant de la somme de 400~g\'en\'erations de la
  r\`egle~\arule{Anneal} \`a \'et\'e r\'ealis\'ee avec \asoft{gnuplot}, elle
  illustre l'int\'er\^et d'une cha\^ine de traitement combinant g\'en\'eration,
  \bialt{r\'eduction}{transformation} et visualisation d'espaces cellulaires. %
  D'autres d'images r\'esultant de manipulations d'espaces g\'en\'er\'es par la
  r\`egle~\arule{Anneal} sont pr\'esent\'ees dans les figures~%
  \mygetdump{anneal-3d}, %
  \mygetdump{anneal-sum}, %
  \mygetdump{anneal-time-xor}, %
  \mygetplot{anneal-pop+anneal-pop-detail}, %
  \mygetdump{anneal-sum-1k-8k}, %
  \mygetdump{anneal-sum-1k-8k-xv}, %
  \mygetdump{anneal-g128-s11} et en \S{}~\mygetref{about-anneal}}

\mydumpfull[t]{anneal-sum}%
{Visualisation par pixmaps des donn\'ees de la figure~\protect\mygetplot{anneal-sum}}%
{D'autres  vues sur la r\'eduction de la
  figure~\mygetplot{anneal-sum} par pixmaps utilisant des colormaps
  diff\'erentes}%
{Le processus d'\'erosion \'evoqu\'e par cette visualisation~(et celle de la
  figure~\mygetplot{anneal-sum}) est r\'ev\'el\'e par une r\'eduction int\'egrale.}

\mydumpfull[t]{life-time-xor}%
{Visualisation par bitmaps de s\'equences de r\'eductions dif\-f\'erentiel\-les
  de la r\`egle~\arule{Life}}%
{Six g\'en\'erations de la r\`egle~\arule{Life} et leurs r\'eductions dif\-f\'e\-ren\-tiel\-les
  d'ordre~1 et~2}%
{Le r\'esultat de la diff\'erence bit \`a bit se lit de gauche \`a droite sur
  les diagonales descendantes, l'espace~2 de la ligne~2 r\'esulte de la
  diff\'erence de l'espace~1 de la ligne~1 et de l'espace~2 de la
  ligne~1\footnote{%
    Ou l'inverse, le \xor\ \'etant commutatif};
  l'espace~(3,3) de~$(1,2)-(3,1)$, mais aussi de~$(2,2)-(2,3)$.}

%%
%% OBJECTS/TIME/STRUCT
%%
\subsection{Structure du temps:~un axe d'organisation privil\'egi\'e}

Le temps est un vecteur d'espaces de taille fixe associ\'e \`a divers
compteurs et pointeurs. %
La construction des op\'erateurs complexes sur le temps repose sur ces
\'el\'ements, devant permettre un acc\`es plus ou moins structur\'e \`a la
m\'emoire temps. %

Lorsque le syst\`eme est r\'ealis\'e sur une machine s\'equentielle, le temps
fournit les espaces temporaires~(accumulateurs) n\'ecessaires \`a la
r\'ealisation de la fonction d'application synchrone de la fonction de
transition sur l'ensemble des cellules de l'espace courant. %
La notion d'accumulateur recevant le r\'esultat de l'application
d'op\'erateurs sur les espaces est sans doute n\'ecessaire dans tous les
contextes, mais sa r\'ealisation sur une machine parall\`ele arbitraire,
ou d\'edi\'ee AC pourrait s'av\'erer d\'elicate\footnote{%
  Sur une architecture SIMD privil\'egiant le nombre de PE~(Processeur
  \'El\'ementaire) au d\'etriment de la m\'emoire locale \`a chaque PE~(de type
  GAPP), o\`u chaque PE est associ\'e \`a une cellule, et la m\'emoire locale
  de chaque PE aux \'etats de chaque cellule, la source et la cible de
  l'application d'une fonction de transition sont potentiellement la
  m\'emoire totale de la machine.}. %

%%
%% OBJECTS/TIME/EDIT
%%
\subsection{\'Edition du temps:~une m\'etaphore cin\'ematogra\-phique}

Voici les op\'erateurs d'\'edition du temps:
\begin{itemize}
\item Remise \`a z\'ero et incr\'ementation d'un compteur d'application
  d'une fonction de transition\footnote{%
    Ce compteur doit \^etre ensuite associ\'e \`a un espace. %
    En r\`egle g\'en\'erale l'inclusion d'un espace dans un objet de plus
    haut niveau implique la cr\'eation de champs nouveaux servant \`a
    m\'emoriser le contexte d'usage de cet espace. %
    Par exemple, la photo d'un espace doit si possible contenir des
    annotations sur les conditions exp\'erimentales de sa cr\'eation.}. %
\item Positionnement absolu et relatif du point d'espace courant. %
\item Test du point courant, relatif aux bornes. %
\item Acc\`es au point d'espace pass\'e. %
\item Application d'une fonction de transition de l'espace courant sur
  l'accumulateur. %
\item \'Echange du point d'espace courant et de l'accumulateur. %
\end{itemize}

Ces primitives permettent la construction de proc\'edures composites. %
Par exemple, pour enregistrer une s\'equence d'espaces d'indice de temps
multiple de~$n$, il faut, apr\`es initialisation de l'espace d'indice
z\'ero et d'un \aquote{timer} \`a~$n$:
\begin{itemize}
\item Appliquer la fonction transition du point courant sur
  l'accumulateur,
\item D\'ecr\'ementer le timer, et si z\'ero incr\'ementer le point courant,
\item \'Echanger le point courant et l'accumulateur. %
  Optionnellement, estampiller le point courant avec le compteur
  d'application de la fonction de transition,
\item It\'erer tant que le point courant n'est pas le dernier,
\end{itemize}

Ces proc\'edures doivent pouvoir \^etre construites par l'utilisateur
aussi simplement que possible, par l'interm\'ediaire d'un interpr\'eteur. %

%%
%% OBJECTS/TIME/REDUCT
%%
\subsection{R\'eductions du temps:~sommes et diff\'erences}

Il existe plusieurs types d'op\'eration de r\'eductions sur le temps. %
Les r\'eductions des s\'equences de r\'eductions d'espaces constituent un
deuxi\`eme niveau de r\'eduction arithm\'etique et statistique. %
On peut par exemple calculer la dimension fractale de l'\'evolution de
n'importe quelle quantit\'e associ\'ee \`a un espace. %

Les r\'eductions d'information int\'egrale\footnote{%
  La relation entre les notions classiques de d\'eriv\'ee et d'int\'egrale et
  le contenu informationnel nous pose probl\`eme. %
  Dans notre contexte de production d'objets discrets et finis, %
  les d\'eriv\'es et les int\'egrales ne sont pas des valeurs limites
  mais des cha\^ines de bits. %
  Le contenu informationnel est la taille de cette cha\^ine de bit, il
  ne se confond pas avec la quantit\'e d'information. %
  Une op\'eration d'int\'egration discr\`ete d\'efinie comme une somme sur une
  direction de l'espace temps cellulaire produit un espace cellulaire
  dont la taille en bits est inf\'erieure \`a la taille de l'espace
  initial. %
  Par extension, nous consid\'erons toute op\'eration de r\'eduction du
  contenu informationnel par diminution du nombre de dimensions de
  l'espace cellulaire comme une int\'egration. %
  
  Une transformation diff\'erentielle~(dont le prototype est le \xor~(xor) bit
  \`a bit ou la distance de Hamming cha\^ine \`a cha\^ine) ne r\'eduit pas le
  contenu informationnel, mais le but poursuivi est a priori une
  r\'eduction de la quantit\'e d'information, ce qui nous am\`ene \`a parler
  d'op\'eration de r\'eduction diff\'erentielle.}%
sur les espaces permettent de passer d'un ensemble d'espaces \`a un
espace unique\footnote{%
  La taille en bit des cellules de l'espace r\'esultat peut \^etre
  sup\'erieure \`a celle des cellules des espaces arguments.}. %
Par exemple un espace dont chaque cellule est la somme
$\sum_{i=0}^{i=n}t_{i}xy$ des valeurs de l'ensemble des cellules de
coordonn\'ees spatiales~$xy$ identiques, sur une s\'equence~$t$ de $n$~espaces. %

\noindent%
La figure~\mygetplot{anneal-sum} montre une vue perspective d'un
espace somme de 400~g\'en\'e\-rations de la r\`egle~\arule{Anneal}. %
L'espace est de taille \asquare{64}. %
L'espace initial a une densit\'e al\'eatoire\,\undemi. %

\noindent%
La figure~\mygetdump{anneal-sum} montre quatre visions pixmap
du m\^eme espace utilisant 4~greymap diff\'erentes. %

Une autre r\'eduction simple est le \xor\ plan de bits \`a plan de bits de
deux espaces, qui identifie les points d'espace modifi\'es d'un point de
temps \`a son successeur. %
Appliqu\'e \`a un AC tel que le jeu de la vie qui produit des patterns
statiques ou de cycle~2 en nombres, cet
op\'erateur permet d'\'eliminer les points d'espace %
de cycle~1~($t_{n-1} \xor t_{n}$),
ou~1 et~2~($t_{n-2} \xor t_{n}$). %

\noindent%
La premi\`ere ligne de la figure~\mygetdump{life-time-xor} montre
6~g\'en\'erations de~\arule{Life}~(de~100 \`a~106) pour un espace de
\asquare{64} cellules et un espace initial al\'eatoire de
densit\'e\,\undemi, la deuxi\`eme ligne le r\'esultat de l'\'elimination de
cellules de cycle~1 et la troisi\`eme de l'\'elimination des cycles~1 et~2. %
Bien s\^ur, ces \'eliminations ne laissent pas intactes les zones non \'elimin\'ees. %

\myplotfull{lifex}%
{\'Evolution de populations de s\'equences de r\'eductions dif\-f\'erentiel\-les}%
{\'Evolution d'une trajectoire compl\`ete de~\arule{Life}}%
{Population de~\arule{Life} et Cycles~1 et~2 sur~(1280 g\'en\'erations)}

\noindent%
La figure~\mygetplot{lifex} montre l'\'evolution sur 1280~g\'en\'erations
des populations associ\'ees aux cycles~1 et~2, pour un espace de
\asquare{128} cellules~(densit\'e al\'eatoire initiale\,\undemi)\footnote{%
  Dans ce qui suit nous utilisons le logarithme pour normaliser les
  r\'esultats}. %

\noindent%
$\text{life} = \log(P_0)$, %
$\text{cycle}^1 = \log(P_0 - P_1)$, %
$\text{cycle}^2 = \log(|P_1 - P_2|)$ %
et $\text{cycle}^{> 2} = \log (P2)$. %

\noindent%
Avec $P_0$ pour la population standard, %
$P_1$ pour la population moins les points d'espace %
de $\text{cycle}^1 = t_{n-1} \xor t_{n}$, %
et $P2$ pour la population moins les points d'espace %
de $\text{cycle}^{<3} = t_{n-2} \xor t_{n}$). %

\noindent%
A la g\'en\'eration~1280, il ne reste plus que des cellules
engag\'ees dans un cycle~1 ou~2. %

\mytwoplotfull{lifex-sum-1}{lifex-sum-2}%
{Une r\'eduction int\'egrale du temps et une r\'eduction diff\'erentiel\-le d'espace}%
{Somme de 512~g\'en\'erations de~\arule{Life} et~\arule{Life} filtr\'e}%
{Une r\'eduction int\'egrale du temps sur une s\'equence de 512~g\'en\'erations
  de~\arule{Life}~(\`a gauche), et sur la s\'equence obtenu par r\'eduction
  diff\'erentielle d'espace avec l'op\'erateur \xor. %
  Cette comparaison sugg\`ere l'existence pour~\arule{Life} d'un
  invariant d'\'echelle pour la densit\'e d'occupation de l'espace
  cellulaire, apr\`es filtrage. %
  Nous consid\'erons cette question, ainsi que la relation entre cette question et les
  exp\'eriences de~\cite{Bak93} sur l'auto-organisation critique
  de~\arule{Life}\footnote{%
    Celle-ci portent sur la taille des avalanches~(le nombre
    d'\'ev\`enements cellulaires \'els\'ementaire, \ie de bit flip) r\'esultant
    d'une perturbation al\'eatoire d'un cycle limite.}, %
  comme une source d'exp\'eriences futures.}

\mydumpfull{life-3d-New}%
{Une visualisation 3D du temps, et d'un temps de r\'eduction diff\'erentiel\-le d'espace}%
{Une autre vue de la figure~\mygetplot{lifex-sum-1+lifex-sum-2}}%
{Une visualisation~3D du temps de~\arule{Life} en partie gauche, %
  et d'un temps de r\'eduction diff\'erentielle
  d'espace de~\arule{Life}~\mygetplot[paren]{lifex-sum-1+lifex-sum-2}.}

\noindent%
La figure~\mygetplot{lifex-sum-1+lifex-sum-2} montre les contours de
la somme de 512~g\'en\'erations de~\arule{Life}~(\`a gauche)
et~\arule{Life} moins les points d'espace de cycle~1 et~2~(\`a droite). %

\mytable{lifex-sum-1}%
{R\'eductions statistiques sur un espace obtenu par r\'eduction int\'egrale
  du temps,}%
{\mysoftcite{pgmtexture} sur \mygetplot{lifex-sum-1+lifex-sum-2}~(gauche)}%
{Caract\'eristiques texturales de l'espace de la
  figure~\mygetplot{lifex-sum-1+lifex-sum-2}~(version non
  filtr\'ee):~exemple de r\'eductions statistiques sur un espace obtenu par
  r\'eduction int\'egrale du temps.}%

\mytable{lifex-sum-2}%
{et par r\'eduction int\'egrale d'une s\'equence
  temporelle de r\'e\-duc\-tion diff\'erentiel\-le d'espaces.}%
{\mysoftcite{pgmtexture} sur \mygetplot{lifex-sum-1+lifex-sum-2}~(droite)}%
{Caract\'eristiques texturales de l'espace de la
  figure~\protect\mygetplot{lifex-sum-1+lifex-sum-2}~(version
  filtr\'ee):~\`a partir d'une r\'eduction int\'egrale d'une s\'equence
  temporelle de r\'eduction diff\'erentielle d'espaces.}

\noindent Les tableaux~\mygettbl{lifex-sum-1}
et~\mygettbl{lifex-sum-2} montrent le r\'esultat de la commande
\mysoftcite{pgmtexture} sur deux espaces somme de la
figure~\mygetplot{lifex-sum-1+lifex-sum-2}. %

La r\'ealisation d'une proc\'edure de mesure de la proportion des cellules
actives engag\'ees dans un point fixe ou un cycle de longueur~2 ou~3, au
top level du Workbench, par combinaison d'objets et de primitives du
noyau ou de la librairie d'application, constitue un objectif
exemplaire de la souplesse d'usage recherch\'ee. %

\mydumpfull{anneal-time-xor}%
{Visualisation par bitmaps de s\'equences de r\'eductions dif\-f\'erentiel\-les
  de la r\`egle~\arule{Anneal}}%
{Six g\'en\'erations de la r\`egle~\arule{Anneal} et leurs r\'eductions dif\-f\'e\-ren\-tiel\-les
  d'ordre~1 et~2}%
{Le mode de production de cette s\'equence est identique \`a celui de la
  figure~\mygetdump{life-time-xor}, son observation\footnote{%
    Plus exactement, l'observation interactive de nombreuses s\'equences
    pour une taille d'espace et un nombre de g\'en\'erations plus
    important.} %
    sugg\`ere comme pour~\arule{Life}, la recherche d'invariants
    d'\'echelle\footnote{%
      \mygetref[phrase]{about-anneal}}}

\mytwoplotfull[t]{anneal-pop}{anneal-pop-detail}%
{\'Evolution de la population de la r\`egle~\arule{Anneal}}%
{La population de 5640~g\'en\'erations de la r\`egle~\arule{Anneal}}%
{Le plot de gauche est un agrandissement sur le cycle limite atteint par
  cette trajectoire de anneal. %
  La d\'ecouverte de l'existence de cycles limite de grande
  taille\footnote{%
    Du m\^eme ordre de grandeur que l'ar\^ete de l'espace, ici~127. %
    } %
  fut une surprise\footnote{%
    Ni la litt\'erature mentionnant la r\`egle:~%
    \arule{M46789} pour~\cite{Vichniac86}, %
    \arule{Lava Lamp}~(une version~3D) pour~\cite{Pickover92b}, %
    \arule{Anneal} pour~\cite{Toffoli87}; %
    ni une exp\'erimentation restreinte ne permettent d'attendre autre
    chose qu'un point fixe ou un cycle limite court~(\ie d'une taille
    ind\'ependante de la taille de l'espace).}. %
  \mygetref[phrase]{about-anneal}}

\mytwoplotfull[t]{heart-pop}{heart-pop-detail}%
{\'Evolution des deux plans de population de la r\`egle~\arule{Heart}}
{512~g\'en\'erations d'un espace de \asquare{256} de la r\`egle~\arule{Heart}}
{La population du plan z\'ero est en abscisse, la population du plan un
  en ordonn\'ee. %
  Cette figure illustre la recherche de possibles attracteurs de
  faibles dimensions, extrait de syst\`emes dot\'es de tr\`es grand nombre de
  degr\'es de libert\'es. %
  }

\noindent%
La figure~\mygetdump{anneal-time-xor} montre le m\^eme traitement
appliqu\'e aux g\'en\'erations~2 \`a~12, par pas de~2, sur la fonction anneal
dans les m\^emes conditions d'espace. %

%% reduction diff\'erentielle du temps sur des r\'eductions int\'egrale
%% d'espace
%%
Nous avons exp\'eriment\'e des r\'eductions diff\'erentielles du temps sur des
s\'equences de la r\'eduction int\'egrale \'el\'ementaire d'espace que constitue
le nombre de bits d'un plan cellulaire. %
Ces r\'eductions sont d\'efinies par une diff\'erence locale des
populations: %

\newcommand{\pop}{\operatorname{Pop}}
\newcommand{\spac}{\operatorname{Space}}

\[\pop_{i}=\pop(\spac_{i})-\pop(\spac_{i-1})\]

%%
%% s\'equence de reduction int\'egrale~(d'espace obtenus par r\'eductions diff\'erentielle
%% d'espace)
%%
Nous cherchons \`a \'etablir la relation entre de telles s\'equences et des
s\'equences de r\'eductions int\'egrales~(population) d'espaces; %
ces espaces \'etant obtenus par la r\'eduction diff\'erentielle spatiale
\'el\'ementaire r\'esultant de l'application du \xor\ plan \`a plan:

\[\pop_{i}=\pop(\spac_{i}\xor\spac_{i-1})\]

Les r\'eductions diff\'erentielles associatives sont r\'ealisables par
it\'eration cumul\'ee d'un op\'erateur \`a nombre d'arguments fixe sur une
fen\^etre de temps. %
On peut alors les consid\'erer comme l'application d'une fonction de
transition sur un espace composite. %

Par exemple, comment transformer une s\'equence de plan de bits d'un
espace~(ou copier cette transformation en conservant la s\'equence
initiale), en s\'equence des diff\'erences plan \`a plan en utilisant
l'op\'erateur \xor. %
\mysetref{time-xor} La fonction de transition \arule{Xor} a deux
arguments d'un bit, et un bit de r\'esultat. %
Ses param\`etres formels sont li\'es aux variables de connexion
\aquote{past} et \aquote{present}~(au lieu de \aquote{east} et
\aquote{west} par exemple), et son r\'esultat est li\'e \`a la variable de
connexion \aquote{present}~(au lieu de \aquote{center}). %
Lors de l'application de la fonction de transition l'espace r\'esultant
est stock\'e dans l'accumulateur. %
Un s\'equenceur peut alors soit copier l'accumulateur dans le point
d'espace courant d'un temps accumulateur de la transformation du temps
courant, soit dans le point d'espace pass\'e du temps courant. %

Plus g\'en\'eralement, si l'on consid\`ere le temps comme un espace de
dimension sup\'erieure, on peut r\'ealiser des projections, ou extraire
des coupes de cet espace de dimension sup\'erieure sur un espace de
dimension inf\'erieure. %
Par exemple, une ligne d'espace de dimension~2 permet d'extraire une
tranche~(un espace de dimension~2) d'un temps d'espace de dimension~2. %

%%
%% OBJECTS/TIME/TRANSFORM
%%
\subsection{Transformations du temps:~l'espace-temps}

Les transformations sur le temps r\'esultent soit d'une combinaison
\bialt{extraction}{construction} qui permet de r\'ealiser des films de
coupes, soit de l'application sur chaque point d'espace du temps d'une
transformation d'espace, soit de l'application d'une transformation
d'espace sur un temps consid\'er\'e comme un espace de dimension
sup\'erieure. %

%%
%% OBJECTS/TIME/VISUAL
%%
\subsection{Visualisation du temps:~s\'equences de r\'eductions et
  projections de plongements}

Les visualisations du temps utilisent les techniques \'evoqu\'ees \`a propos
de la visualisation des espaces. %
Celles-ci sont pilot\'ees par un s\'equenceur qui combine les primitives
de manipulation du temps et les primitives de visualisation de
l'espace. %
Les m\^emes outils sont utilis\'es pour visualiser les espaces r\'esultant
de r\'eduction du temps. %

La visualisation des s\'equences de r\'eductions d'espace utilise des
plotters pour pr\'esenter des graphes \bialt{temps}{variable}, %
ou multi-variables~(population plan~1, population plan~2 e.g.), %
ou toutes combinaisons de variables, de r\'eductions et d'expressions
autoris\'ees par les outils off-line utilisables. %

\noindent%
La figure~\mygetplot{anneal-pop+anneal-pop-detail}~(anneal) pr\'esente
l'\'evolution de la population de la r\`egle~\arule{Anneal} pendant
5640~g\'en\'erations \'echantillonn\'ees chaque 20~g\'en\'erations. %
La taille de l'espace est \asquare{127}. %
On peut y voir un cycle limite. %

\noindent%
La figure~\mygetplot{anneal-pop+anneal-pop-detail}~(anneal d\'etail)
pr\'esente un agrandissement sur le cycle. %
La d\'etection et la mesure de cycles limites sur les populations sont des
exemples de r\'eductions appliqu\'ees sur une s\'equence de r\'eduction. %

\noindent%
La figure~\mygetplot{heart-pop+heart-pop-detail}~(heart) pr\'esente
l'\'evolution de la population du plan~0 contre l'\'evolution de la
population du plan~1 de la r\`egle \arule{Heart} pendant 512~g\'en\'erations d'un
espace de \asquare{256}. %

\noindent%
On peut y voir le d\'ecouplage des deux variables. %
La figure~\mygetplot{heart-pop+heart-pop-detail}~(heart d\'etail) est un
agrandissement sur le d\'ecouplage. 

\clearpage

%%
%% Local Variables:
%% mode: latex
%% mode: auto-fill
%% iso-accents-minor-mode: nil
%% TeX-master: "top"
%% TeX-command-default: "mytex"
%% Local IspellParsing: "latex-mode ~latin1"
%% Local IspellPersDict: "~/.ispell_lisp"
%% LocalWords: "s\'equenceur SIMD PE GAPP incr\'ementation timer D\'ecr\'ementer It\'erer"
%% LocalWords: "incr\'ementer Optionnellement interpr\'eteur anneal-sum anneal xor"
%% LocalWords: "plot-anneal-sum fractale pixmap greymap CA patterns lifex.ps"
%% LocalWords: "life-time-xor width height lifex-sum lifex-sum texturales level"
%% LocalWords: "pgmtexture Workbench anneal-time-xor past east west center multi"
%% LocalWords: "plotters off-line anneal-pop anneal-pop-detail heart-pop"
%% LocalWords: "heart-pop-detail versus heart informationnel Hamming gnuplot"
%% LocalWords: "k-xv anneal-g about-anneal bitmaps lifex Pop Space time-xor e.g"
%% End:
